\appendix

\section{Proof Details}\label{sec:proof-details}

\subsection{Rounded canonical oracle and transcript}\label{app:step-001}

\begin{lemma}[Endpoint-grid rounding bound]\label{lem:step-001-grid-rounding}
Under Assumption~\ref{assump:sq-parameter-regime}, let
\[
K=\lceil 1/\tau\rceil,
\qquad
G=\{-1+2j/K:0\le j\le K\},
\]
and let \(\rho:[-1,1]\to G\) be the fixed deterministic nearest-grid
map.  Then, for every \(u\in[-1,1]\),
\[
|\rho(u)-u|\le \frac1K\le\tau.
\]
Moreover, \(\rho(-1)=-1\), \(\rho(1)=1\), either fixed choice at an
interval midpoint has error exactly \(1/K\), and the same bound holds when
\(\tau\ge 1\).
\end{lemma}

\begin{proof}
Since \(\tau>0\),
\[
K=\lceil 1/\tau\rceil\ge 1,
\qquad
K\ge \frac1\tau,
\qquad
\frac1K\le\tau.
\tag{A.1}\label{eq:grid-inversion}
\]
For \(0\le j\le K\), put \(\gamma_j=-1+2j/K\).  Then
\[
\gamma_0=-1,
\qquad
\gamma_K=1,
\qquad
\gamma_{j+1}-\gamma_j=\frac2K
\quad(0\le j<K).
\tag{A.2}\label{eq:grid-spacing}
\]
The intervals \([\gamma_j,\gamma_{j+1}]\) cover \([-1,1]\).  If
\(u\in[\gamma_j,\gamma_{j+1}]\), then
\[
\min\{u-\gamma_j,\gamma_{j+1}-u\}
\le \frac{\gamma_{j+1}-\gamma_j}{2}
=\frac1K.
\tag{A.3}\label{eq:grid-half-spacing}
\]
Every other grid point is no closer to \(u\) than one of these adjacent
endpoints.  The nearest-grid definition and
\eqref{eq:grid-inversion}--\eqref{eq:grid-half-spacing} therefore give the
claimed chain.

At either endpoint the nearest-grid error is zero.  At a midpoint both
adjacent points are at distance \(1/K\), so the fixed tie rule preserves the
bound and makes \(\rho\) single valued.  If \(\tau\ge1\), then
\(0<1/\tau\le1\), hence \(K=1\), \(G=\{-1,1\}\), and the covering radius
is \(1=1/K\le\tau\), including the midpoint \(u=0\).
\end{proof}

\begin{proposition}[Validity of the canonical oracle]
\label{prop:step-001-canonical-oracle}
Under Assumptions~\ref{assump:sq-parameter-regime}
and~\ref{assump:universal-adversarial-sq} and
Lemma~\ref{lem:step-001-grid-rounding}, for every distribution
\(\mathcal D\) on \(\mathcal X\) and every \(h\in\mathcal H\), the rule
\[
\mathcal O^\rho_{\mathcal D,h}(q)
=\rho\!\left(\mathbb E_{x\sim\mathcal D}q(x,h(x))\right)
\]
is tolerance-valid for every current unrestricted SQ query, including a
query chosen adaptively from the preceding transcript.
\end{proposition}

\begin{proof}
Fix \((\mathcal D,h)\).  For an unrestricted query
\(q:\mathcal X\times\{+1,-1\}\to[-1,1]\), pointwise boundedness gives
\[
-1\le u_q:=\mathbb E_{x\sim\mathcal D}q(x,h(x))\le1.
\]
Consequently, Lemma~\ref{lem:step-001-grid-rounding} gives
\[
\left|
\mathcal O^\rho_{\mathcal D,h}(q)
-\mathbb E_{x\sim\mathcal D}q(x,h(x))
\right|
=|\rho(u_q)-u_q|
\le\frac1K\le\tau.
\tag{A.4}\label{eq:querywise-validity}
\]
The calculation holds for every bounded query, and hence at every query
selected from any reached adaptive prefix.  Validity is querywise: the
discrepancies in \eqref{eq:querywise-validity} are not summed over rounds.
Endpoint expectations are rounded exactly, and the fixed midpoint rule gives
a unique reply of error \(1/K\).  This is precisely the required adaptive
absolute-tolerance condition.
\end{proof}

\begin{proposition}[Unique canonical padded path]
\label{prop:step-001-padded-path}
Under Assumptions~\ref{assump:sq-parameter-regime}
and~\ref{assump:universal-adversarial-sq} and
Proposition~\ref{prop:step-001-canonical-oracle}, for every
\((\mathcal D,h,r)\), the interaction of \(A_r\) with
\(\mathcal O^\rho_{\mathcal D,h}\) determines a unique actual reply prefix
of length at most \(m\).  Appending \(-1\in G\) after stopping determines a
unique \(z^{\mathcal D,h,r}\in G^m\), and the appended entries are not
oracle replies.  If \(m=0\), this path is the unique empty string.
\end{proposition}

\begin{proof}
Fix \((\mathcal D,h,r)\).  The canonical oracle is a deterministic rule
with values in \(G\).  At the empty prefix, the taped protocol \(A_r\)
either stops or issues one uniquely determined first query.  Inductively,
after the actual replies \((z_1,\ldots,z_{t-1})\) have been fixed, taped
determinism gives one stopping decision or one next query \(q_t\).  In the
latter case the next reply is uniquely
\[
z_t=\rho\!\left(
\mathbb E_{x\sim\mathcal D}q_t(x,h(x))
\right)\in G.
\tag{A.5}\label{eq:canonical-recursion}
\]
The query budget makes this recursion terminate with a unique actual prefix
\((z_1,\ldots,z_s)\), where \(0\le s\le m\).  If \(s<m\), define
\[
z_{s+1}=\cdots=z_m=-1.
\tag{A.6}\label{eq:fixed-padding}
\]
Here \(-1=-1+2\cdot0/K\in G\).  Since \eqref{eq:fixed-padding} is appended
only after stopping, none of these entries is supplied to the oracle or
requires a tolerance check.  The deterministic recursion, midpoint rule,
and padding value yield a unique element of \(G^m\).

When \(m=0\), no query is permitted, no reply or padding entry is produced,
and the zero-fold product is \(G^0=\{\emptyset\}\).  Thus the path is
\(\emptyset\) and oracle validity along it is vacuous.
\end{proof}

\begin{proof}[Completion of the oracle and transcript construction]
Lemma~\ref{lem:step-001-grid-rounding} establishes the rounding inequality
at endpoints, midpoint ties, and \(\tau\ge1\).  Proposition
\ref{prop:step-001-canonical-oracle} applies it to every current adaptive
query, and Proposition~\ref{prop:step-001-padded-path} constructs the unique
full path, including \(\emptyset\) at \(m=0\).  No \(m/K\) error is
introduced because only executed replies are checked and the post-stop
suffix is never executed.
\end{proof}

\subsection{Finite response-tree map and measurable law}
\label{app:step-002}

\begin{proposition}[Total padded response tree]
\label{prop:step-002-total-replay}
Under Assumptions~\ref{assump:sq-parameter-regime}
and~\ref{assump:universal-adversarial-sq} and the grid of
Lemma~\ref{lem:step-001-grid-rounding}, every complete tape \(r\) and every
\(z=(z_1,\ldots,z_m)\in G^m\) determine by prescribed-reply replay a
terminal binary predictor
\[
g_{r,z}:\mathcal X\to\{+1,-1\}.
\]
If replay stops after reading exactly \(s<m\) replies, then every
\(\widetilde z\in G^m\) satisfying
\(\widetilde z_t=z_t\) for \(1\le t\le s\) obeys
\[
g_{r,\widetilde z}=g_{r,z}
\quad\text{as functions on }\mathcal X.
\]
Thus suffix symbols are unread, the terminal output can be copied to every
depth-\(m\) descendant, and the fixed all-\(-1\) suffix selects one such
continuation.  If \(m=0\), \(g_{r,\emptyset}\) is the taped learner's
no-query terminal predictor.
\end{proposition}

\begin{proof}
Fix \(r\) and \(z\in G^m\).  Complete-tape conditioning makes the protocol
deterministic.  At the empty prefix it either stops or issues a uniquely
defined first query.  More generally, after the entries
\((z_1,\ldots,z_{t-1})\) actually supplied, the next query, stopping
decision, and terminal output are deterministic functions of that prefix.
If query \(t\) is issued, replay supplies \(z_t\).  Since \(K>0\),
\[
-1\le -1+\frac{2j}{K}\le1
\qquad(0\le j\le K),
\]
so all prescribed replies lie in \([-1,1]\).  At most \(m\) queries can be
executed, after which the protocol returns a binary predictor.  This proves
existence and totality of every \(g_{r,z}\).  No instance, population
expectation, or tolerance certificate is needed for this synthetic replay.

Suppose replay on \(z\) stops after \(s<m\) entries and \(\widetilde z\)
has the same first \(s\) entries.  Inductively for \(t=0,\ldots,s\), the
two replays have the same state after \(t\) replies: the initial taped state
is common, and equality of states and of the next prescribed reply gives the
same deterministic transition.  After \(s\) replies both executions make
the same stopping decision and return the same predictor.  Neither requests
entry \(s+1\), so later entries cannot affect the output.  This includes
root stopping \(s=0\); when \(s=m\), there is no suffix to check.

The padded tree assigns the same output to every continuation of a terminal
prefix.  The value \(-1\in G\) selects one continuation but is not read.
When \(m=0\), no response is supplied and the root output is exactly
\(g_{r,\emptyset}\).
\end{proof}

\begin{lemma}[Exact response-string count]
\label{lem:step-002-response-count}
Under Assumption~\ref{assump:sq-parameter-regime} and
Lemma~\ref{lem:step-001-grid-rounding}, the grid values are pairwise
distinct and
\[
|G|=K+1,
\qquad
|G^m|=(K+1)^m
=\bigl(\lceil1/\tau\rceil+1\bigr)^m.
\]
For \(m=0\), \(G^0=\{\emptyset\}\) and both sides equal one.
\end{lemma}

\begin{proof}
The grid lemma gives \(K\ge1\).  If \(0\le j<j'\le K\), then
\[
\left(-1+\frac{2j'}K\right)-\left(-1+\frac{2j}K\right)
=\frac{2(j'-j)}K>0.
\tag{A.7}\label{eq:grid-distinct}
\]
Thus \(|G|=K+1\).  By definition,
\(|G^0|=1=(K+1)^0\).  If \(|G^t|=(K+1)^t\), unique decomposition into a
length-\(t\) prefix and a last grid value gives
\[
|G^{t+1}|=|G^t|\,|G|=(K+1)^t(K+1)=(K+1)^{t+1}.
\tag{A.8}\label{eq:product-count-induction}
\]
Finite induction and \(K=\lceil1/\tau\rceil\) prove the claim.
\end{proof}

\begin{proposition}[Finite response-tree feature map]
\label{prop:step-002-feature-map}
Under Assumptions~\ref{assump:sq-parameter-regime}
and~\ref{assump:universal-adversarial-sq},
Proposition~\ref{prop:step-002-total-replay}, and
Lemma~\ref{lem:step-002-response-count}, fix the lexicographic order on
\(G^m\) and put \(N=|G^m|=(K+1)^m\).  Then every complete tape \(r\)
defines the total map
\[
\phi_r(x)=\bigl(g_{r,z}(x)\bigr)_{z\in G^m}:
\mathcal X\to\{+1,-1\}^{N}\subseteq\mathbb R^N.
\]
This construction is valid on arbitrary, including infinite, domains, after
stopping at any prefix, and when distinct response strings yield identical
terminal predictors.  It uses only the fixed learner and tape, \(m,\tau\),
the grid and its order, and the padding convention, so it is independent of
\((\mathcal D,h)\).  At \(m=0\),
\(\phi_r(x)=(g_{r,\emptyset}(x))\in\mathbb R\).
\end{proposition}

\begin{proof}
Proposition~\ref{prop:step-002-total-replay} supplies one binary function
\(g_{r,z}\) at every element of the common finite index set \(G^m\), and
Lemma~\ref{lem:step-002-response-count} gives exactly \(N\) indices.  Their
fixed order forms an \(N\)-tuple, all of whose values at each \(x\) belong
to \(\{+1,-1\}\).

Every coordinate comes from replay of a fixed taped protocol on a fixed
numeric string.  Replay never chooses an instance, evaluates a population
query, or follows a realized canonical path, so \(\phi_r\) is fixed before
\((\mathcal D,h)\).  The construction evaluates a finite list of total
functions and uses no cardinality or topological property of \(\mathcal X\).
Suffix invariance keeps every coordinate defined after early stopping.  If
\(g_{r,z}=g_{r,z'}\) for \(z\ne z'\), the equal functions simply occupy two
separately labeled coordinates; the dimension counts response strings, not
distinct function values.  At \(m=0\), the index set is
\(\{\emptyset\}\), so the displayed scalar map follows.
\end{proof}

\begin{proposition}[Pre-instance pushforward feature law]
\label{prop:step-002-preinstance-law}
Under Assumptions~\ref{assump:sq-parameter-regime}
and~\ref{assump:universal-adversarial-sq} and
Proposition~\ref{prop:step-002-feature-map}, let
\((\Omega_R,\Sigma_R,\mu)\) be the complete-tape probability space, with
\(R\) the identity map.  Define
\[
\mathfrak F_N=\{\Phi:\mathcal X\to\mathbb R^N\},
\qquad
T(r)=\phi_r,
\]
and equip \(\mathfrak F_N\) with
\[
\mathscr A_T
=\{B\subseteq\mathfrak F_N:T^{-1}(B)\in\Sigma_R\}.
\tag{A.9}\label{eq:final-sigma-algebra}
\]
Then \(\mathscr A_T\) is a sigma algebra, \(T\) is measurable, and
\[
\mathsf P_A=T_{\#}\mu,
\qquad
\mathsf P_A(B)=\mu(T^{-1}(B))
\quad(B\in\mathscr A_T),
\tag{A.10}\label{eq:pushforward-law}
\]
is a probability law on maps into \(\mathbb R^N\), where
\(N=(\lceil1/\tau\rceil+1)^m\).  The measurable space and law are fixed
before and shared by every \((\mathcal D,h)\).

For each fixed \((\mathcal D,h)\), coordinate evaluations, the functional
\[
F_{\mathcal D,h}(\Phi)
=\inf_{w\in\mathbb R^N}
L_{\mathrm{tie}}\bigl(\mathcal D,h,
x\mapsto\langle w,\Phi(x)\rangle\bigr),
\tag{A.11}\label{eq:optimal-map-loss}
\]
and the tape-selected basis loss are measurable.  Moreover,
\[
\begin{aligned}
\mathbb E_{\Phi\sim\mathsf P_A}F_{\mathcal D,h}(\Phi)
&=\mathbb E_{R\sim\mu}F_{\mathcal D,h}(\phi_R)\\
&\le \mathbb E_{R\sim\mu}
L_{\mathrm{tie}}\!\left(
\mathcal D,h,
x\mapsto\left\langle
e_{z^{\mathcal D,h,R}},\phi_R(x)
\right\rangle\right).
\end{aligned}
\tag{A.12}\label{eq:measurable-candidate-interface}
\]
Distinct tapes may induce the same map; no inverse-tape selection is needed.
\end{proposition}

\begin{proof}
Let \(\mathscr X\) be the measurable structure on \(\mathcal X\) used by
the distributions, queries, predictors, and losses.  The notation for a
randomized learner and its complete tape entails measurable finite
execution: for every fixed \(z\in G^m\),
\[
(r,x)\longmapsto g_{r,z}(x)
\tag{A.13}\label{eq:joint-replay-measurability}
\]
is \((\Sigma_R\otimes\mathscr X)\)-measurable.  This is the measurable
output typing needed for the setting's displayed tape expectations, not an
additional statistical premise.  Since \(G^m\) is finite,
\[
(r,x)\longmapsto T(r)(x)
=\bigl(g_{r,z}(x)\bigr)_{z\in G^m}
\tag{A.14}\label{eq:joint-map-measurability}
\]
is jointly measurable into \(\{+1,-1\}^N\).

We first verify the parameter integral used below.  For a
product-measurable \(Q\subseteq\Omega_R\times\mathcal X\), write
\(Q_r=\{x:(r,x)\in Q\}\), and define
\[
\mathcal L_{\mathcal D}
=\left\{Q\in\Sigma_R\otimes\mathscr X:
r\mapsto\mathcal D(Q_r)\text{ is }\Sigma_R\text{-measurable}\right\}.
\]
Every section \(Q_r\) is in \(\mathscr X\), because the sets with this
property form a sigma algebra containing the measurable rectangles.  The
class \(\mathcal L_{\mathcal D}\) is a Dynkin system: the full-space
section measure is one; complements replace \(\mathcal D(Q_r)\) by
\(1-\mathcal D(Q_r)\); and for pairwise disjoint \(Q_j\),
\[
\mathcal D\!\left(\left(\bigcup_{j\ge1}Q_j\right)_r\right)
=\sum_{j\ge1}\mathcal D((Q_j)_r),
\tag{A.15}\label{eq:section-disjoint-union}
\]
the pointwise limit of measurable partial sums.  For a measurable rectangle
\(E\times C\),
\[
\mathcal D((E\times C)_r)=\mathbf 1_E(r)\mathcal D(C),
\]
which is measurable.  The measurable rectangles are a pi-system generating
\(\Sigma_R\otimes\mathscr X\); Dynkin's pi-lambda theorem
\citep[Section~1]{Kallenberg2021} therefore gives
\(\Sigma_R\otimes\mathscr X\subseteq\mathcal L_{\mathcal D}\).

If \(a:\Omega_R\times\mathcal X\to[0,1]\) is jointly measurable, choose
nonnegative product-measurable simple functions \(a_n\uparrow a\).  The set
calculation above and finite linearity make
\(r\mapsto\int a_n(r,x)\,d\mathcal D(x)\) measurable, and monotone
convergence gives
\[
\int a_n(r,x)\,d\mathcal D(x)
\uparrow
\int a(r,x)\,d\mathcal D(x).
\tag{A.16}\label{eq:parameter-integral-limit}
\]
Thus
\[
r\longmapsto\int a(r,x)\,d\mathcal D(x)
\tag{A.17}\label{eq:parameter-integral}
\]
is measurable.

We next establish the feature-map measurable space.  Equation
\eqref{eq:final-sigma-algebra} contains \(\mathfrak F_N\), is closed under
complements because
\[
T^{-1}(\mathfrak F_N\setminus B)
=\Omega_R\setminus T^{-1}(B),
\]
and is closed under countable unions because
\[
T^{-1}\!\left(\bigcup_{j\ge1}B_j\right)
=\bigcup_{j\ge1}T^{-1}(B_j).
\]
It is therefore a sigma algebra, and \(T\) is measurable by definition.
Equation \eqref{eq:pushforward-law} defines a probability measure: inverse
images preserve disjoint unions and
\[
\mathsf P_A(\mathfrak F_N)
=\mu(T^{-1}(\mathfrak F_N))=\mu(\Omega_R)=1.
\]
For every nonnegative \(\mathscr A_T\)-measurable \(H\), indicators,
finite linearity, and increasing simple approximation prove
\[
\int_{\mathfrak F_N}H(\Phi)\,d\mathsf P_A(\Phi)
=\int_{\Omega_R}H(T(r))\,d\mu(r).
\tag{A.18}\label{eq:pushforward-integral}
\]

The final sigma algebra has the exact factorization property
\[
U:\mathfrak F_N\to\mathbb R\text{ is }\mathscr A_T\text{-measurable}
\quad\Longleftrightarrow\quad
U\circ T\text{ is }\Sigma_R\text{-measurable}.
\tag{A.19}\label{eq:factorization-measurability}
\]
Indeed, for each Borel \(C\subseteq\mathbb R\),
\[
T^{-1}(U^{-1}(C))=(U\circ T)^{-1}(C),
\]
and the two directions follow directly from
\eqref{eq:final-sigma-algebra}.  In particular, a measurable tape statistic
constant on fibers of \(T\) factors uniquely on \(T(\Omega_R)\) through a
measurable map statistic.

For \(x\in\mathcal X\) and \(z\in G^m\), let
\(\operatorname{ev}_{x,z}(\Phi)=\Phi_z(x)\).  Its pullback is
\(g_{r,z}(x)\), measurable by \eqref{eq:joint-replay-measurability}, so
\eqref{eq:factorization-measurability} makes every coordinate evaluation
measurable.  For fixed \(w\in\mathbb R^N\), define
\[
\ell_{\mathcal D,h,w}(\Phi)
=L_{\mathrm{tie}}\bigl(
\mathcal D,h,x\mapsto\langle w,\Phi(x)\rangle\bigr).
\tag{A.20}\label{eq:fixed-weight-loss}
\]
By \eqref{eq:joint-map-measurability}, the function
\[
(r,x)\longmapsto
\mathbf 1\{\langle w,T(r)(x)\rangle h(x)\le0\}
\]
is jointly measurable and bounded.  Equation
\eqref{eq:parameter-integral} and then
\eqref{eq:factorization-measurability} show that every fixed-weight loss,
including each standard-basis loss, is measurable.

An uncountable infimum of measurable functions need not be measurable, so
we reduce it without changing ties.  Let
\(\mathcal V_N=\{+1,-1\}^N\).  Declare \(w,w'\in\mathbb R^N\) equivalent
when
\[
\mathbf 1\{\langle w,v\rangle y\le0\}
=\mathbf 1\{\langle w',v\rangle y\le0\}
\quad\text{for every }(v,y)\in\mathcal V_N\times\{+1,-1\}.
\tag{A.21}\label{eq:tie-pattern-equivalence}
\]
There are finitely many classes because this is a binary pattern on a finite
set.  Choose one representative of every nonempty class and adjoin all
standard basis vectors \(e_z\), obtaining a finite
\(W_N\subseteq\mathbb R^N\).  Since every map in \(T(\Omega_R)\) is
binary-valued, equivalent weights have the same tie-error indicator at every
point, including at zero inner products.  Hence
\[
F_{\mathcal D,h}(T(r))
=\min_{w\in W_N}\ell_{\mathcal D,h,w}(T(r)).
\tag{A.22}\label{eq:finite-pattern-minimum}
\]
This finite minimum is measurable on the tape space, and
\eqref{eq:factorization-measurability} proves measurability of the exact
functional in \eqref{eq:optimal-map-loss}.  No topology, separability, or
finiteness of \(\mathcal X\) is used.

For the selected-basis interface, put
\(Z_{\mathcal D,h}(r)=z^{\mathcal D,h,r}\in G^m\).  Its measurability
follows by finite induction through measurable execution: the taped
next-query and stopping maps are measurable,
\eqref{eq:parameter-integral} makes each current population query mean
measurable in \(r\), \(\rho\) is Borel because it is piecewise constant on
finitely many intervals with fixed midpoint choices, and post-stop padding
is fixed.  Define
\[
C_{\mathcal D,h}(r)
=\ell_{\mathcal D,h,e_{Z_{\mathcal D,h}(r)}}(T(r)).
\]
Since \(G^m\) is finite,
\[
C_{\mathcal D,h}(r)
=\sum_{z\in G^m}
\mathbf 1\{Z_{\mathcal D,h}(r)=z\}
\ell_{\mathcal D,h,e_z}(T(r)),
\tag{A.23}\label{eq:selected-basis-finite-sum}
\]
which is measurable.  For every \(r\), the definition of the infimum gives
\[
F_{\mathcal D,h}(T(r))\le C_{\mathcal D,h}(r).
\tag{A.24}\label{eq:tape-infimum-candidate}
\]
Apply \eqref{eq:pushforward-integral} to \(F_{\mathcal D,h}\), then
integrate \eqref{eq:tape-infimum-candidate}, to obtain
\eqref{eq:measurable-candidate-interface}.

If \(T(r)=T(r')\) but the two canonical indices differ, the same mapwise
infimum is separately bounded by the two tape-side candidates; no inverse of
\(T\) or common coordinate on a fiber is used.  All objects in
\eqref{eq:final-sigma-algebra}--\eqref{eq:pushforward-law} depend only on the
learner, tape law, finite grid, ordering, and padding.  Verifying the
functionals after fixing \((\mathcal D,h)\) does not alter the sigma algebra
or law.  At \(m=0\), \(T(r)\) has one root coordinate,
\(Z_{\mathcal D,h}\equiv\emptyset\), and all formulas reduce to their
one-coordinate forms.
\end{proof}

\begin{proof}[Completion of the response-tree construction]
Proposition~\ref{prop:step-002-total-replay} constructs every binary replay
predictor and proves exact suffix invariance.  Lemma
\ref{lem:step-002-response-count} gives
\[
|G^m|=(K+1)^m=(\lceil1/\tau\rceil+1)^m=N.
\]
Proposition~\ref{prop:step-002-feature-map} forms the total map on the common
index set, retaining equal predictors as legal separate coordinates.
Proposition~\ref{prop:step-002-preinstance-law} then constructs the
measurable pushforward, the finite tie-pattern reduction, and the tape-side
candidate comparison.  These conclusions place the maps under one law fixed
before the instance and cover root and intermediate stopping, duplicate map
fibers, arbitrary domains, and the empty-product case.
\end{proof}

\subsection{Exact replay and tie-free coordinate}\label{app:step-003}

\begin{lemma}[Prefix synchronization]
\label{lem:step-003-prefix-synchronization}
Under Assumption~\ref{assump:universal-adversarial-sq} and
Propositions~\ref{prop:step-001-padded-path}
and~\ref{prop:step-002-total-replay}, fix \((\mathcal D,h,r)\), put
\(z=z^{\mathcal D,h,r}\), and let
\(s\in\{0,\ldots,m\}\) be the number of replies read before \(A_r\)
terminates against \(\mathcal O^\rho_{\mathcal D,h}\).  After every
\(t\in\{0,\ldots,s\}\) replies, the actual interaction and the replay
indexed by \(z\) have read the same prefix and occupy the
same taped configuration.  If \(t<s\), they issue the same next query and
receive the same next reply; at \(t=s\), their stopping decision and
terminal binary predictor coincide.  This includes \(s=0\) and \(m=0\).
\end{lemma}

\begin{proof}
Induct on \(t\).  At \(t=0\), both executions use the same tape, protocol,
and empty history, so they have the same initial configuration and first
action.  If this action is terminal, then \(s=0\) and the outputs coincide.

Suppose \(0\le t<s\) and the configurations and prefixes agree after
\(t\) replies.  Taped determinism gives the same next query
\(q_{t+1}\).  The actual reply is
\[
\mathcal O^\rho_{\mathcal D,h}(q_{t+1})
=\rho\!\left(
\mathbb E_{x\sim\mathcal D}q_{t+1}(x,h(x))
\right),
\tag{A.25}\label{eq:synchronized-reply}
\]
which, by the canonical-path recursion, equals
\(z^{\mathcal D,h,r}_{t+1}\).  Prescribed-string replay receives that same
entry.  The two executions therefore apply the same transition to the same
state and reply, proving agreement after \(t+1\) replies.

At \(t=s\), equal configurations force the same next action; because the
actual run stops, replay stops and returns the same predictor.  If \(s=m\),
the query budget forces termination no later than this state.  If \(m=0\),
then \(s=0\) and the base case applies.  The induction uses the identical
canonical reply, rather than two merely tolerance-valid replies, so midpoint
choices, \(\tau\ge1\), and arbitrary adaptive dependence create no defect.
\end{proof}

\begin{proposition}[Canonical terminal replay]
\label{prop:step-003-terminal-replay}
Under Assumption~\ref{assump:universal-adversarial-sq},
Lemma~\ref{lem:step-003-prefix-synchronization}, and
Proposition~\ref{prop:step-002-total-replay}, for every
\((\mathcal D,h,r)\),
\[
g_{r,z^{\mathcal D,h,r}}
=A_r^{\mathcal O^\rho_{\mathcal D,h}}(\mathcal D,h)
\tag{A.26}\label{eq:terminal-replay-identity}
\]
Both sides are functions from \(\mathcal X\) to \(\{+1,-1\}\).
If stopping occurs after \(s<m\) queries, the padded suffix is unread and
every full continuation of the same terminal prefix gives the same output.
At \(m=0\), this is equality of the no-query root predictors.
\end{proposition}

\begin{proof}
Lemma~\ref{lem:step-003-prefix-synchronization} puts both executions in the
same terminal configuration after the same \(s\) replies, so their terminal
predictor objects agree and hence agree at every \(x\).  When \(s<m\), the
canonical suffix is appended only after actual stopping; synchronized replay
also stops at that prefix and requests none of those entries.  Proposition
\ref{prop:step-002-total-replay} shows that all other continuations have the
same output.  This includes root stopping.  For \(m=0\), the sole response
index is \(\emptyset\) and the synchronization base case gives the root
identity.
\end{proof}

\begin{proposition}[One-hot selection of the actual binary output]
\label{prop:step-003-one-hot-score}
Under Assumption~\ref{assump:universal-adversarial-sq} and
Propositions~\ref{prop:step-002-feature-map}
and~\ref{prop:step-003-terminal-replay}, for every
\((\mathcal D,h,r,x)\),
\[
\begin{aligned}
\left\langle e_{z^{\mathcal D,h,r}},\phi_r(x)\right\rangle
&=g_{r,z^{\mathcal D,h,r}}(x)\\
&=A_r^{\mathcal O^\rho_{\mathcal D,h}}(\mathcal D,h)(x)
\in\{+1,-1\}.
\end{aligned}
\tag{A.27}\label{eq:one-hot-output}
\]
The identity requires neither distinct coordinate functions nor an injective
tape-to-map construction.  For \(m=0\), it uses
\(e_{\emptyset}=(1)\in\mathbb R\).
\end{proposition}

\begin{proof}
The feature map has coordinates
\(\phi_r(x)=(g_{r,z'}(x))_{z'\in G^m}\).  The standard basis vector at
\(z^{\mathcal D,h,r}\) therefore gives by direct multiplication
\[
\left\langle e_{z^{\mathcal D,h,r}},\phi_r(x)\right\rangle
=g_{r,z^{\mathcal D,h,r}}(x).
\tag{A.28}\label{eq:coordinate-selection}
\]
Proposition~\ref{prop:step-003-terminal-replay} supplies equality with the
actual output, and both terminal predictors are binary.  Repeated functions
remain separately indexed, and equal maps from different tapes satisfy
\eqref{eq:coordinate-selection} separately on each tape.  At \(m=0\), the
one-coordinate formula is multiplication by \((1)\).  The basis vector is
chosen after fixing \((\mathcal D,h,r)\), and does not alter the previously
constructed map or law.
\end{proof}

\begin{proposition}[Exact tie-free loss transfer]
\label{prop:step-003-tie-free-transfer}
Under Assumption~\ref{assump:universal-adversarial-sq} and
Propositions~\ref{prop:step-002-preinstance-law}
and~\ref{prop:step-003-one-hot-score}, for every fixed
\((\mathcal D,h)\), tape \(r\), and point \(x\),
\[
\begin{aligned}
\mathbf 1\!\left\{
\left\langle e_{z^{\mathcal D,h,r}},\phi_r(x)\right\rangle h(x)\le0
\right\}
&=\mathbf 1\{g_{r,z^{\mathcal D,h,r}}(x)h(x)<0\}\\
&=\mathbf 1\!\left\{
A_r^{\mathcal O^\rho_{\mathcal D,h}}(\mathcal D,h)(x)h(x)<0
\right\}.
\end{aligned}
\tag{A.29}\label{eq:tie-event-identity}
\]
Consequently,
\[
\begin{aligned}
&L_{\mathrm{tie}}\!\left(
\mathcal D,h,
x\mapsto\left\langle e_{z^{\mathcal D,h,r}},\phi_r(x)\right\rangle
\right)\\
&\quad=L_{\mathrm{tie}}(\mathcal D,h,g_{r,z^{\mathcal D,h,r}})
=L_{\mathrm{bin}}(\mathcal D,h,g_{r,z^{\mathcal D,h,r}})\\
&\quad=L_{\mathrm{bin}}\!\left(
\mathcal D,h,A_r^{\mathcal O^\rho_{\mathcal D,h}}(\mathcal D,h)
\right).
\end{aligned}
\tag{A.30}\label{eq:tie-loss-identity}
\]
The left side is a measurable tape function.  No selector from a feature-map
fiber to a tape or basis vector is required.
\end{proposition}

\begin{proof}
By Proposition~\ref{prop:step-003-one-hot-score}, the selected score, replay
output, and actual output are the same sign.  Since \(h(x)\) is also a sign,
their product belongs to \(\{+1,-1\}\), so zero is impossible and
\[
\{\langle e_{z^{\mathcal D,h,r}},\phi_r(x)\rangle h(x)\le0\}
=\{\langle e_{z^{\mathcal D,h,r}},\phi_r(x)\rangle h(x)=-1\}
=\{\langle e_{z^{\mathcal D,h,r}},\phi_r(x)\rangle h(x)<0\}.
\tag{A.31}\label{eq:no-tie-event}
\]
Substitute the pointwise output identities to obtain
\eqref{eq:tie-event-identity}; taking \(\mathcal D\)-probabilities gives
\eqref{eq:tie-loss-identity}.  The nonzero score is a consequence of binary
coordinate selection, not a margin assumption.

Proposition~\ref{prop:step-002-preinstance-law} proves measurability of the
finite-valued path \(Z_{\mathcal D,h}(r)=z^{\mathcal D,h,r}\) and of
\[
r\longmapsto L_{\mathrm{tie}}\!\left(
\mathcal D,h,
x\mapsto\left\langle e_{Z_{\mathcal D,h}(r)},T(r)(x)\right\rangle
\right).
\tag{A.32}\label{eq:measurable-selected-loss}
\]
Since \(T(r)=\phi_r\), equation \eqref{eq:tie-loss-identity} identifies
\eqref{eq:measurable-selected-loss} with the actual binary loss, which is
therefore measurable.  If equal feature maps arise from tapes with different
canonical indices, the comparison remains tape-wise; no common basis vector
or inverse of \(T\) is asserted.  Repeated coordinates and \(m=0\) are
covered by the same binary identity.
\end{proof}

\begin{proof}[Completion of the replay and loss transfer]
Lemma~\ref{lem:step-003-prefix-synchronization} proves that adaptive replay
has zero prefix defect.  Proposition~\ref{prop:step-003-terminal-replay}
makes the padded suffix unread and identifies the two terminal predictors.
Proposition~\ref{prop:step-003-one-hot-score} then identifies one
standard-basis coordinate with the actual learner output, and Proposition
\ref{prop:step-003-tie-free-transfer} converts the pointwise binary identity
into exact tie- and binary-loss equality.  These conclusions remain
tape-side on duplicate map fibers and reduce to the root coordinate when
\(m=0\).
\end{proof}

\subsection{Expected-risk and dimension certificate}\label{app:step-004}

\begin{proposition}[Tape-wise infimum comparison]
\label{prop:step-004-infimum-pushforward}
Under Assumptions~\ref{assump:sq-parameter-regime}
and~\ref{assump:universal-adversarial-sq} and
Propositions~\ref{prop:step-001-padded-path},
\ref{prop:step-002-feature-map}, and
\ref{prop:step-002-preinstance-law}, fix \((\mathcal D,h)\) and let
\(F_{\mathcal D,h}\) be defined by
\eqref{eq:optimal-map-loss}.  Then
\[
\begin{aligned}
\mathbb E_{\Phi\sim\mathsf P_A}F_{\mathcal D,h}(\Phi)
&=\mathbb E_{R\sim\mu}F_{\mathcal D,h}(\phi_R)\\
&\le\mathbb E_{R\sim\mu}
L_{\mathrm{tie}}\!\left(
\mathcal D,h,
x\mapsto\left\langle
e_{z^{\mathcal D,h,R}},\phi_R(x)
\right\rangle\right).
\end{aligned}
\tag{A.33}\label{eq:infimum-pushforward}
\]
The comparison is tape-wise and remains valid on noninjective fibers of
\(r\mapsto\phi_r\).
\end{proposition}

\begin{proof}
For fixed \(r\), the canonical path gives
\(z^{\mathcal D,h,r}\in G^m\), so the feature-map coordinate construction
makes \(e_{z^{\mathcal D,h,r}}\in\mathbb R^N\) an admissible weight.
Therefore
\[
F_{\mathcal D,h}(\phi_r)
\le L_{\mathrm{tie}}\!\left(
\mathcal D,h,
x\mapsto\left\langle e_{z^{\mathcal D,h,r}},\phi_r(x)\right\rangle
\right).
\tag{A.34}\label{eq:pointwise-infimum-bound}
\]
No minimizer is claimed.  Proposition~\ref{prop:step-002-preinstance-law}
gives measurability, boundedness in \([0,1]\), and the pushforward identity.
Integrating \eqref{eq:pointwise-infimum-bound} and using monotonicity yields
\eqref{eq:infimum-pushforward}.

If \(T(r)=T(r')\) but the canonical indices differ, the common left side of
\eqref{eq:pointwise-infimum-bound} is separately bounded by each legal
tape-side candidate.  The expectation retains those candidates on their
original tapes, so no common fiberwise weight, inverse map, or measurable
selection is used.
\end{proof}

\begin{proposition}[Risk closure at the canonical oracle]
\label{prop:step-004-canonical-risk}
Under Assumptions~\ref{assump:sq-parameter-regime}
and~\ref{assump:universal-adversarial-sq} and
Propositions~\ref{prop:step-001-canonical-oracle}
and~\ref{prop:step-003-tie-free-transfer}, for every
\((\mathcal D,h)\),
\[
\begin{aligned}
&\mathbb E_{R\sim\mu}
L_{\mathrm{tie}}\!\left(
\mathcal D,h,
x\mapsto\left\langle e_{z^{\mathcal D,h,R}},\phi_R(x)\right\rangle
\right)\\
&\quad=\mathbb E_{R\sim\mu}
L_{\mathrm{tie}}(\mathcal D,h,g_{R,z^{\mathcal D,h,R}})\\
&\quad=\mathbb E_{R\sim\mu}
L_{\mathrm{bin}}\!\left(
\mathcal D,h,
A_R^{\mathcal O^\rho_{\mathcal D,h}}(\mathcal D,h)
\right)
\le\varepsilon.
\end{aligned}
\tag{A.35}\label{eq:canonical-risk-closure}
\]
The outer randomness is only the learner tape, and rounding contributes no
tie-risk remainder.
\end{proposition}

\begin{proof}
Fix \((\mathcal D,h)\).  Proposition
\ref{prop:step-001-canonical-oracle} makes
\(\mathcal O^\rho_{\mathcal D,h}\) tolerance-valid.  The universal
adversarial-SQ guarantee applies to this particular valid policy and gives
\[
\mathbb E_{R\sim\mu}
L_{\mathrm{bin}}\!\left(
\mathcal D,h,
A_R^{\mathcal O^\rho_{\mathcal D,h}}(\mathcal D,h)
\right)
\le\varepsilon.
\tag{A.36}\label{eq:primitive-canonical-bound}
\]
This is an instantiation of an every-valid-policy premise, not a favorable
oracle assumption.  Proposition~\ref{prop:step-003-tie-free-transfer} gives
the two preceding equalities in \eqref{eq:canonical-risk-closure} for each
tape and supplies measurability.  Integrating only over \(R\sim\mu\) and
using \eqref{eq:primitive-canonical-bound} proves the result.  The population
loss retains its defining integration over \(x\sim\mathcal D\); no instance,
oracle, or additional random variable is averaged externally.
\end{proof}

\begin{proposition}[Simultaneous tie-risk and dimension certificate]
\label{prop:step-004-dimension-certificate}
Under Assumptions~\ref{assump:sq-parameter-regime}
and~\ref{assump:universal-adversarial-sq},
Lemma~\ref{lem:step-002-response-count}, and
Propositions~\ref{prop:step-002-preinstance-law},
\ref{prop:step-004-infimum-pushforward}, and
\ref{prop:step-004-canonical-risk}, the single pre-instance law
\(\mathsf P_A\) satisfies
\[
\sup_{\mathcal D}\sup_{h\in\mathcal H}
\mathbb E_{\Phi\sim\mathsf P_A}
\left[
\inf_{w\in\mathbb R^N}
L_{\mathrm{tie}}\bigl(
\mathcal D,h,x\mapsto\langle w,\Phi(x)\rangle
\bigr)
\right]
\le\varepsilon.
\tag{A.37}\label{eq:uniform-risk-certificate}
\]
Consequently,
\[
\operatorname{dc}^{\mathrm{tie}}_\varepsilon(\mathcal H)
\le N=\bigl(\lceil1/\tau\rceil+1\bigr)^m.
\tag{A.38}\label{eq:dimension-certificate}
\]
At \(m=0\), \(N=1\); at \(\varepsilon=0\), the expected tie risk is
zero; and at \(\tau\ge1\), \(N=2^m\).
\end{proposition}

\begin{proof}
Fix arbitrary \((\mathcal D,h)\).  Propositions
\ref{prop:step-004-infimum-pushforward}
and~\ref{prop:step-004-canonical-risk} give
\[
\begin{aligned}
&\mathbb E_{\Phi\sim\mathsf P_A}
\left[
\inf_{w\in\mathbb R^N}
L_{\mathrm{tie}}\bigl(
\mathcal D,h,x\mapsto\langle w,\Phi(x)\rangle
\bigr)
\right]\\
&\quad\le\mathbb E_{R\sim\mu}
L_{\mathrm{tie}}\!\left(
\mathcal D,h,
x\mapsto\left\langle e_{z^{\mathcal D,h,R}},\phi_R(x)\right\rangle
\right)\\
&\quad=\mathbb E_{R\sim\mu}
L_{\mathrm{tie}}(\mathcal D,h,g_{R,z^{\mathcal D,h,R}})\\
&\quad=\mathbb E_{R\sim\mu}
L_{\mathrm{bin}}\!\left(
\mathcal D,h,
A_R^{\mathcal O^\rho_{\mathcal D,h}}(\mathcal D,h)
\right)
\le\varepsilon.
\end{aligned}
\tag{A.39}\label{eq:complete-risk-chain}
\]
The first relation evaluates the mapwise infimum; the middle relations are
exact loss identities; and the last inequality is the primitive guarantee
at a valid oracle.  There is no additive term.

Proposition~\ref{prop:step-002-preinstance-law} constructs the same
\(\mathsf P_A\), measurable map space, and ambient dimension before the
arbitrary pair is chosen.  Taking the two suprema in
\eqref{eq:complete-risk-chain} therefore gives
\eqref{eq:uniform-risk-certificate}, with no union bound or probability-mode
change.  Lemma~\ref{lem:step-002-response-count} gives exactly
\[
N=|G^m|=(K+1)^m
=\bigl(\lceil1/\tau\rceil+1\bigr)^m.
\tag{A.40}\label{eq:complete-dimension-count}
\]
The definition of \(\operatorname{dc}^{\mathrm{tie}}_\varepsilon\) asks
for one pre-instance law with this uniform expected mapwise loss property;
\eqref{eq:uniform-risk-certificate} and
\eqref{eq:complete-dimension-count} exhibit it.

When \(m=0\), \(G^0=\{\emptyset\}\), \(N=1\), and the same chain uses the
no-query root predictor.  When \(\varepsilon=0\), nonnegativity and
\eqref{eq:uniform-risk-certificate} make the expected tie loss exactly zero,
with no rounding remainder.  If \(\tau\ge1\), then \(K=1\), so the exact
count is \(2^m\).  The total-map and pushforward results already cover
arbitrary domains and noninjective tape-to-map fibers.
\end{proof}

\begin{proof}[Completion of the expected-risk and dimension argument]
For an arbitrary \((\mathcal D,h)\), Proposition
\ref{prop:step-004-infimum-pushforward} first transfers the mapwise
optimal-loss functional through the pre-instance pushforward and evaluates
the infimum at the tape-specific canonical basis vector.  Proposition
\ref{prop:step-004-canonical-risk} then uses the exact tie-free identity and
the every-valid-oracle premise to bound this candidate loss by
\(\varepsilon\), with expectation only over the tape.  Proposition
\ref{prop:step-004-dimension-certificate} observes that the same law was
fixed before the arbitrary instance, takes the two suprema, and applies the
definition with the exact count
\(N=(\lceil1/\tau\rceil+1)^m\).  No term is absorbed, and the
\(m=0\), \(\varepsilon=0\), and \(\tau\ge1\) reductions are exact.
\end{proof}

\subsection{Proof of the Main Theorem}\label{app:main-proof}

\begin{proof}[Proof of Theorem~\ref{thm:main}]
Lemma~\ref{lem:step-001-grid-rounding} and
Propositions~\ref{prop:step-001-canonical-oracle}
and~\ref{prop:step-001-padded-path} produce a legal canonical policy and its
path.  Proposition~\ref{prop:step-002-total-replay},
Lemma~\ref{lem:step-002-response-count}, and
Propositions~\ref{prop:step-002-feature-map}
and~\ref{prop:step-002-preinstance-law} construct the finite response-tree
maps and their one pre-instance law.  Lemma
\ref{lem:step-003-prefix-synchronization} and
Propositions~\ref{prop:step-003-terminal-replay},
\ref{prop:step-003-one-hot-score}, and
\ref{prop:step-003-tie-free-transfer} identify a canonical binary coordinate
with the actual learner output and preserve the tie/binary baseline exactly.
Finally, Propositions~\ref{prop:step-004-infimum-pushforward}
and~\ref{prop:step-004-canonical-risk} give the exact comparison chain in
Theorem~\ref{thm:main}, while
Proposition~\ref{prop:step-004-dimension-certificate} takes the universal
instance suprema and applies the dimension definition.  These conclusions
are precisely its risk and dimension statements, including the
\(m=0\), \(\varepsilon=0\), and \(\tau\ge1\) boundary cases.
\end{proof}
